\section{Influence Domains}
\label{sec:influence-domains}

The cost model in Section~\ref{sec:task-cost-model} defines attempt-attributed
cost: active resource demand and worker occupancy. A complete task-behavior
model, however, also contains variables that are not cost outputs: task
identity, intrinsic work, execution context, symptoms, outcomes, diagnostics,
risks, and profiling-policy triggers. These variables must not be collapsed into
one metric space, because they play different roles in estimation, diagnosis,
risk interpretation, and profiling-policy selection.

Influence domains assign semantic roles to model variables. They are not
telemetry sources and not metric categories. A metric, span event, log entry,
queue event, profile sample, or domain event becomes part of the model only
after it is mapped to a model variable with an assigned influence domain. The
same raw observation may support more than one model variable, but each graph
node has one primary role.

\subsection{Domain Set}
\label{sec:influence-domain-set}

Let
\[
\mathcal{D}
=
\{
D_{\mathrm{id}},
D_{\mathrm{intr}},
D_{\mathrm{ctx}},
D_{\mathrm{res}},
D_{\mathrm{occ}},
D_{\mathrm{sym}},
D_{\mathrm{out}},
D_{\mathrm{diag}},
D_{\mathrm{risk}},
D_{\mathrm{policy}}
\}
\]
be the set of influence domains. The domains represent task identity, intrinsic
features, execution context, resource demand, worker occupancy, symptoms,
outcomes, reliability diagnostics, operational risk, and profiling-policy
triggers.

The domains classify roles rather than data sources. For example, a queue event
may be mapped to an outcome node, a worker lifecycle event may be mapped to an
occupancy node, and a log entry may be mapped to a symptom node. The mapping
depends on how the resulting variable is used in the task-behavior model.

\subsection{Domain Assignment}
\label{sec:influence-domain-assignment}

For a Task Behavior Graph associated with task kind $k$, let $V_k$ be the node
set. The domain assignment is a typing function
\[
\delta_k : V_k \rightarrow \mathcal{D}.
\]
When the task kind is clear from context, we write $\delta(v)$ instead of
$\delta_k(v)$.

This function assigns one primary influence domain to each graph node. If a raw
signal supports several roles, the mapping layer should create several model
variables rather than assigning several roles to one node. For example, a
dependency-timeout log entry may support a symptom variable, an attempt-outcome
variable, and risk evidence, but those roles should be represented by distinct
nodes:
\[
\delta(\texttt{dependency\_timeout\_symptom}) = D_{\mathrm{sym}},
\]
\[
\delta(\texttt{attempt\_timeout\_outcome}) = D_{\mathrm{out}},
\]
\[
\delta(\texttt{retry\_storm\_risk}) = D_{\mathrm{risk}}.
\]
This one-domain rule makes graph interpretation and validation explicit.

\subsection{Domain Semantics}
\label{sec:influence-domain-semantics}

Table~\ref{tab:influence-domain-semantics} summarizes the domains, their roles,
typical examples, and common confusions. The table is part of the model
definition: it specifies how variables may be interpreted before they are
connected by graph edges.

\begingroup
\renewcommand{\arraystretch}{1.18}
\setlength{\tabcolsep}{4pt}

\begin{longtable}{
p{0.12\linewidth}
p{0.26\linewidth}
p{0.28\linewidth}
p{0.24\linewidth}
}
\caption{Influence domains and their semantic roles.}
\label{tab:influence-domain-semantics}\\

\toprule
Domain & Role & Examples & Common confusion \\
\midrule
\endfirsthead

\toprule
Domain & Role & Examples & Common confusion \\
\midrule
\endhead

\bottomrule
\endfoot

$D_{\mathrm{id}}$ &
Identifies which task is being modeled and selects reusable task-kind
structure. &
task kind, task version, queue, tenant, producer, priority, deadline &
Identity selects a model; it is not a cost value. \\

$D_{\mathrm{intr}}$ &
Describes work contained in the task instance. &
payload size, record count, object count, batch size, partition size, operation
mode &
Intrinsic work is not execution context. \\

$D_{\mathrm{ctx}}$ &
Describes the environment and state under which the attempt executes. &
dependency health, network condition, disk pressure, cache state, worker-pool
pressure, retry state &
Context may influence cost, but it is not an attempt-attributed cost output. \\

$D_{\mathrm{res}}$ &
Describes active resource demand or observed active resource usage attributed to
the attempt. &
CPU time, memory peak, allocation volume, disk I/O, network I/O, accelerator
time &
Resource demand is not bounded-capacity occupancy. \\

$D_{\mathrm{occ}}$ &
Describes bounded execution capacity held by the attempt over time. &
worker-slot time, lease time, connection hold time, lock hold time,
semaphore-permit time &
Occupancy is not wall-clock duration by itself. \\

$D_{\mathrm{sym}}$ &
Describes an intermediate observed signal that indicates relevant or abnormal
behavior. &
dependency wait spike, retry burst, lock-contention signal, queue-delay spike &
A symptom is not necessarily the root cause or final outcome. \\

$D_{\mathrm{out}}$ &
Describes how the attempt completed. &
success, failure, timeout, cancellation, retryable error, dead-lettering, lease
expiration &
An outcome is not a cost value by itself. \\

$D_{\mathrm{diag}}$ &
Describes reliability of estimates or model fit. &
residual magnitude, missingness, low sample count, drift, tail ratio,
multimodality, confidence &
A diagnostic is not an operational consequence by itself. \\

$D_{\mathrm{risk}}$ &
Describes consequence-weighted operational danger. &
queue starvation, retry storm, SLO violation, replication lag, tenant impact,
cascading overload &
Risk is not active resource demand or worker occupancy. \\

$D_{\mathrm{policy}}$ &
Describes a profiling-policy trigger or observation response. &
minimal accounting, detailed accounting, tracing, tail-aware profiling,
diagnostic profiling &
A profiling-policy trigger is not a scheduling or placement decision. \\

\end{longtable}
\endgroup

\subsection{Permitted Domain Flow}
\label{sec:influence-domain-flow}

The domain assignment constrains graph edges. Let
\[
\leadsto_{\mathcal{D}}
\subseteq
\mathcal{D}\times\mathcal{D}
\]
denote the permitted domain-flow relation. If $(u,v)\in E_k$, then the edge is
well typed only when
\[
\delta_k(u)\leadsto_{\mathcal{D}}\delta_k(v).
\]
This condition does not prove causality. It only states that a graph edge uses
the source variable in an admissible role relative to the target variable.

Typical permitted flows are:
\[
D_{\mathrm{id}}
\leadsto_{\mathcal{D}}
D_{\mathrm{intr}},\,
D_{\mathrm{ctx}},\,
D_{\mathrm{res}},\,
D_{\mathrm{occ}},\,
D_{\mathrm{policy}},
\]
because identity can select task-kind structure, relevant features, relevant
cost dimensions, and policy priors.

Intrinsic features and execution context may influence cost estimates:
\[
D_{\mathrm{intr}},\,D_{\mathrm{ctx}}
\leadsto_{\mathcal{D}}
D_{\mathrm{res}},\,D_{\mathrm{occ}}.
\]
For example, payload size may influence network output bytes, while dependency
health may influence worker-slot time or connection hold time.

Context may also influence symptoms and outcomes:
\[
D_{\mathrm{ctx}}
\leadsto_{\mathcal{D}}
D_{\mathrm{sym}},\,D_{\mathrm{out}}.
\]
For example, degraded dependency health may be associated with dependency-wait
symptoms or timeout outcomes.

Cost, symptoms, and outcomes may feed diagnostics or risk interpretation:
\[
D_{\mathrm{res}},\,D_{\mathrm{occ}},\,D_{\mathrm{sym}},\,D_{\mathrm{out}}
\leadsto_{\mathcal{D}}
D_{\mathrm{diag}},\,D_{\mathrm{risk}}.
\]
For example, a large worker-slot residual may feed reliability diagnostics, and
a timeout outcome may feed retry-storm risk.

Diagnostics and risk may feed profiling-policy triggers:
\[
D_{\mathrm{diag}},\,D_{\mathrm{risk}}
\leadsto_{\mathcal{D}}
D_{\mathrm{policy}}.
\]
For example, low confidence together with high queue-starvation risk may trigger
tail-aware profiling.

Some flows are restricted. Outcomes and diagnostics for attempt $a$ may update
task-kind state for future attempts, but they must not be used as pre-execution
input features for the same attempt. Otherwise, the graph would leak
post-execution information into the estimate it is supposed to produce. In
particular, observed values, residuals, and final outcomes for attempt $a$ must
not be folded back into $z_a$ for the same pre-execution estimate.

\subsection{Ambiguous Signals}
\label{sec:influence-domain-ambiguous-signals}

Some operational signals are ambiguous until the model assigns a role. Duration,
retry count, timeout, and queue delay are common examples.

A duration signal may represent different roles. Attributed CPU duration belongs
to $D_{\mathrm{res}}$. Worker-slot time belongs to $D_{\mathrm{occ}}$.
Dependency wait time is often a symptom in $D_{\mathrm{sym}}$, unless a
descriptor explicitly defines it as an occupancy dimension. Elapsed wall-clock
duration is usually an outcome-related observation, not worker occupancy by
itself.

A retry-count signal may also have several roles. The attempt number available
before execution may be execution context. A burst of retries observed during
execution may be a symptom. A retry scheduled after failure may be an outcome. A
high retry count may also be evidence for retry-storm risk. These roles should
not be represented by one overloaded graph node.

This distinction is important because the same raw observation can enter the
model through several typed variables:
\[
s
\mapsto
\{v_1,\ldots,v_m\},
\qquad
\delta_k(v_j)\in\mathcal{D}.
\]
The mapping layer may derive several variables from one raw signal, but the
graph layer keeps their roles separate.

\subsection{Anti-Double-Counting Rule}
\label{sec:influence-domain-anti-double-counting}

The anti-double-counting rule is central: variables that describe the same
incident must not be aggregated as if they were independent additive costs.
Network instability, dependency wait time, retry count, wall-clock duration,
connection hold time, and worker-slot time may all appear during one degraded
attempt. They do not have the same role.

Consider an attempt that runs during packet loss and waits thirty seconds on a
degraded dependency. The incident may produce the following variables:
\[
\delta(\texttt{network\_condition}) = D_{\mathrm{ctx}},
\]
\[
\delta(\texttt{dependency\_wait\_spike}) = D_{\mathrm{sym}},
\]
\[
\delta(\texttt{retry\_count}) \in \{D_{\mathrm{ctx}},D_{\mathrm{sym}},D_{\mathrm{out}}\},
\]
depending on whether retry count is used as pre-execution attempt state, an
observed retry burst, or a completion result,
\[
\delta(\texttt{wall\_clock\_duration}) = D_{\mathrm{out}},
\]
\[
\delta(\texttt{worker\_slot\_time}) = D_{\mathrm{occ}},
\]
and
\[
\delta(\texttt{queue\_starvation\_risk}) = D_{\mathrm{risk}}.
\]

Only some of these variables are cost dimensions. Network condition is context.
Dependency wait is a symptom unless it is explicitly modeled as a held-capacity
dimension. Wall-clock duration is an elapsed outcome. Worker-slot time is
occupancy only for intervals during which the worker slot is held. Queue
starvation is a consequence, not an additional resource-demand dimension.

The model must assign roles before aggregation. This prevents a single
underlying incident from being counted once as context, again as a symptom,
again as duration, again as occupancy, and again as risk. Role assignment does
not remove the relationships between these variables; it makes those
relationships explicit.

\subsection{Boundaries}
\label{sec:influence-domain-boundaries}

Influence domains define how model variables may be interpreted and connected.
They do not define the cost estimator, do not map raw telemetry into canonical
features, do not compute uncertainty or risk, do not select the profiling policy,
and do not prove causality. Those responsibilities belong to later sections.

The role of this section is narrower: it provides the typing discipline used by
Task Behavior Graphs. Section~\ref{sec:task-behavior-graphs} uses the domain
assignment to connect typed variables into task-kind-specific graphs.
Section~\ref{sec:metrics-signals-features} explains how raw observations are
mapped into model variables. Sections~\ref{sec:entropy-confidence-risk} and
\ref{sec:profiling-policy-selection} consume cost estimates, diagnostics, and
risk to reason about reliability and observation policy.
